\section{Иллюстрации результатов}

Ниже приведены основные визуализации, полученные в ходе выполнения лабораторной работы.

\begin{figure}[H]
\centering
\includegraphics[width=0.95\textwidth]{task1_decision_boundary.png}
\caption{Граница решения для линейно разделимой задачи бинарной классификации.}
\label{fig:lab_01_fig_1}
\end{figure}

\begin{figure}[H]
\centering
\includegraphics[width=0.95\textwidth]{task1_convergence.png}
\caption{Сходимость персептрона в задаче бинарной классификации.}
\label{fig:lab_01_fig_2}
\end{figure}

\begin{figure}[H]
\centering
\includegraphics[width=0.95\textwidth]{task2_xor_problem.png}
\caption{Геометрическая постановка задачи XOR на плоскости.}
\label{fig:lab_01_fig_3}
\end{figure}

\begin{figure}[H]
\centering
\includegraphics[width=0.95\textwidth]{task3_decision_map.png}
\caption{Карта решений персептрона из двух нейронов для четырех классов.}
\label{fig:lab_01_fig_4}
\end{figure}
